\subsection{Conditioning at the complete spectral nodes}\label{cs:conditioning}

The canonical representative of a quotient class has degree less than
\(N\). Off the spectral nodes we specify whether we use this
representative or an actual unreduced product. Define the norm
\[
 \normJ{F}=\max\bigl\{|F(c)|+|F'(c)|:c\text{ a double node};
                     \ |F(-2)|\text{ if }N\text{ is odd}\bigr\}.
\]
It is a norm on the full quotient and is submultiplicative because
\(|FG|+|(FG)'|\le(|F|+|F'|)(|G|+|G'|)\) at each double node.
The derivative at an odd simple node is not part of the data.
Put
\begin{equation}\label{cs:weights}
 w_{N,r}=\begin{cases}e_r(-2),&0\le r\le M,\\1,&M<r<N.\end{cases}
\end{equation}
These weights are positive. The following estimates are purely at the
complete main color; no tree or filling-power hypothesis is needed.

\begin{lemma}[Complete-node conditioning]\label{cs:conditioning-lemma}
There are absolute constants \(C,A\) such that
\begin{align}
 |\phi(F)|&\le CN^A\normJ{F},&
 \normJ{R^\epsilon L_j}&\le CN^A,
 &1\le j\le N,\quad\epsilon\in\{0,1,-1\},\label{cs:finite-cost}\\
 \normJ{e_r}&\le CN^A w_{N,r},&&&0\le r<N,\label{cs:basis-norm}\\
 \log w_{N,r}&=Ng(r/N)+O(\log N),&&&0\le r\le M.\label{cs:weight-rate}
\end{align}
For \(r>M\), every value of \(e_r\) at a spectral node is zero,
but its first normals satisfy \(\normJ{e_r}\le N^2/16\).
Every canonical polynomial also satisfies
\begin{equation}\label{cs:ellipse}
 |F(2\cos(\theta+\ii y))|
       \le CN^A e^{N|y|}\normJ{F},\qquad \theta,y\in\mathbb R.
\end{equation}
This last assertion is not made for arbitrary unreduced representatives.
\end{lemma}
\begin{proof}
The polynomial with the distinct complete nodes as simple roots is
\[
 Q_N(z)=\begin{cases}
 2T_{N/2}(z/2),&N\text{ even},\\
 (z+2)\{U_M(z/2)-U_{M-1}(z/2)\},&N=2M+1.
 \end{cases}
\]
Its interval norm and first derivatives are polynomial in \(N\).
At even roots, \(|Q_N'|\ge N/2\). At an odd interior root
\(c=2\cos t\), differentiation gives
\(|Q_N'(c)|=N/(2\sin(t/2))\); at \(-2\) it gives \(N\).
Each Lagrange cardinal is bounded on \([-2,2]\) by
\(\|Q_N'\|_\infty/|Q_N'(c)|\), by the mean-value formula.
Thus interpolation at the distinct nodes has polynomial cost.

For \(\deg F<N\), interpolate its values by \(I\) with
\(\deg I<\deg Q_N\). Write \(F-I=Q_N H\). At each double node,
\[
 H(c)=\frac{F'(c)-I'(c)}{Q_N'(c)}.
\]
The degree of \(H\) is less than the number of double nodes, so these
values determine it. The interpolating polynomial for this second
step is \(Q_N\) at even levels and \(Q_N/(z+2)\) at odd levels.
The latter has polynomial interval norm, and its root derivatives
have magnitude \(|Q_N'(c)|/(c+2)\ge N/8\). The same cardinal
estimate applies. Derivative bounds do not presume good conditioning:
a degree-\(n\) polynomial has Chebyshev coefficients at most twice its
interval norm, and \(\|T_j'(z/2)\|_\infty\le j^2/2\).
This proves polynomial equivalence of the canonical interval norm,
its complete value-and-normal norm, and its Chebyshev coefficient sum.

In particular \([z^d]T_d(z/2)=1/2\), so the exact functional in
\eqref{cs:turn-trace} has polynomial cost. The circle polynomial has
\(|L_j|\le j\) on \([-2,2]\), and the same derivative estimate gives
a polynomial bound for its first derivative. The crossing has unit
values and logarithmic derivative \(O(N^2)\) at every double node
by \eqref{cs:circle-crossing}. Its inverse has the same bounds.
This proves \eqref{cs:finite-cost}. The Chebyshev coefficient bound
and \(|T_k(\cos(\theta+\ii y))|\le e^{k|y|}\) prove
\eqref{cs:ellipse}.

For low \(r\), the values \(e_r(c)\) at distinct nodes are either
zero or nonnegative, and do not exceed \(e_r(-2)\). Its degree is
less than the number of distinct nodes, so the first interpolation
step bounds its interval norm and derivative by \(CN^A e_r(-2)\).
The product at \(-2\) is
\[
 e_r(-2)=\frac{\prod_{a<r}4\cos^2((a+1/2)\pi/N)}{S_r^2}.
\]
Endpoint sum--integral comparison, as in \eqref{cs:factorial}, gives
\(\log e_r(-2)=Ng(r/N)+O(\log N)\) uniformly in \(r\le M\).
The midpoint logarithmic intervals at the upper endpoint contribute
only \(O(\log N)\), including both parities.

For high \(r\), put \(l=N-r\le\lfloor N/2\rfloor\),
\(\theta=\pi/N\), and \(B_k(z)=\prod_{a<k}(c_a-z)\).
Reflection gives \(B_rB_l=(-1)^NP_N\). At a double node indexed by
\(0\le j<\lfloor N/2\rfloor\), the derivative of \(e_r\) vanishes
unless \(j<l\), and otherwise
\begin{equation}\label{cs:high-normal}
 |e_r'(c_j)|=
 \frac{S_{l-1}^2}
 {2\sin((2j+1)\theta)S_{l+j}S_{l-1-j}}.
\end{equation}
Compare quadratic coefficients in the reflected identity, using
\[
 \frac{|P_N''(c_j)|}{2}=\frac{N^2}{4\sin^2((2j+1)\theta)},\qquad
 |B_l'(c_j)|=\frac{S_{l+j}S_{l-1-j}}{2\sin((2j+1)\theta)},
\]
and \(S_rS_{l-1}=N\). The ratio
\(D_j=S_{l+j}S_{l-1-j}/S_{l-1}^2\) starts at \(2\sin(l\theta)\)
and is nondecreasing, since
\[
 \frac{D_{j+1}}{D_j}
 =\frac{\sin((l+j+1)\theta)}{\sin((l-1-j)\theta)}\ge1;
\]
the two angles sum to at most \(\pi\). Consequently
\eqref{cs:high-normal} is at most
\(1/[4\sin((2j+1)\theta)\sin(l\theta)]\le N^2/16\).
The odd simple-node value is zero. This retains every normal and
completes the proof.
\end{proof}

Throughout the estimates, a uniform error \(o(1)\) may absorb powers
of \(N\), a fixed number of sums and the errors of fixed estimates.
It is uniform over all indicated physical colors and retained signs.
A strict exponential loss is reduced before those errors are absorbed.
